% Generated draft: requires mathematical and source review.
\documentclass[11pt]{article}
\usepackage[margin=25mm]{geometry}
\usepackage{fontspec}
\setmainfont{Latin Modern Roman}
\usepackage{amsmath,amssymb,mathtools,unicode-math}
\setmathfont{Latin Modern Math}
\usepackage{enumitem,needspace,xurl,hyperref}
\setlength{\parindent}{0pt}
\setlength{\parskip}{5pt}
\setlength{\emergencystretch}{3em}
\newcommand{\N}{\mathbb{N}}
\newcommand{\Z}{\mathbb{Z}}
\newcommand{\Q}{\mathbb{Q}}
\newcommand{\R}{\mathbb{R}}
\newcommand{\C}{\mathbb{C}}
\newcommand{\F}{\mathbb{F}}
\newcommand{\A}{\mathbb{A}}
\newcommand{\PP}{\mathbb{P}}
\newcommand{\E}{\mathbb{E}}
\newcommand{\eps}{\varepsilon}
\newcommand{\dd}{\,\mathrm{d}}
\begin{document}
\section*{UnsolvedMath ID: 30006389}
\textbf{Source problem code:} catalogue entry
\subsection{Definitions and mathematical statement}
For a positive integer \(n\), let \([n]=\{1,\ldots,n\}\). A partial block design on \([n]\) is a finite family of subsets \(A_1,\ldots,A_m\subseteq[n]\) satisfying \(|A_i\cap A_j|\leq1\) whenever \(i<j\). A pairwise balanced block design on \([n]\) is understood here in the usual sense: every block has at least two elements, and every unordered pair of distinct elements of \([n]\) lies in exactly one block. The supplied record uses the latter term without explicitly defining it; this interpretation requires source review.

The supplied statement attributes the following two questions to Erd\H{o}s.
\begin{enumerate}
\item Is it true that
\[
\begin{gathered}
\forall c\in(0,1)\ \exists C=C(c)\in\R_{>0}\ \forall n,m\in\Z_{\geq1}\ 
\forall A_1,\ldots,A_m\subseteq[n],\\
\left[\bigl(\forall i,\ |A_i|>c\sqrt n\bigr)\ \land\ 
\bigl(\forall i<j,\ |A_i\cap A_j|\leq1\bigr)\right]\\
\Longrightarrow\ 
\exists B\subseteq[n]\ \forall i\in\{1,\ldots,m\},\quad
0<|B\cap A_i|\leq C\,?
\end{gathered}
\]
Thus \(C\) may depend on \(c\), but must be independent of \(n\), \(m\), and the family of blocks. No requirement that every pair of points be covered is imposed in this question.
\item For \(n\geq2\), a finite sequence of integers
\[
n\geq x_1\geq\cdots\geq x_m\geq2,\qquad m\geq1,
\]
is called block-compatible for \(n\) if there exists a pairwise balanced block design \(A_1,\ldots,A_m\) on \([n]\) with \(|A_i|=x_i\) for every \(i\). Let \(\mathcal S_n\) be the set of these sequences, counting each sequence once regardless of the number of designs realizing it. Is it true that
\[
\exists c_0>0\ \exists n_0\in\Z_{\geq2}\ \forall n\in\Z_{\geq n_0},\qquad
|\mathcal S_n|\geq\exp\!\bigl(c_0\sqrt n\log n\bigr)?
\]
Here \(\log\) denotes the natural logarithm, and \(c_0\) is an absolute constant independent of \(n\), unrelated to the parameter \(c\) in the first question.
\end{enumerate}

These are the two questions in the supplied cleaned statement. The supplied dated literature assessment says that the separate counting lower bound has been settled, while a necessary-and-sufficient characterization of block-compatible sequences remains unresolved. That characterization is not explicitly posed in the supplied mathematical statement. Consequently, the second question is retained as a source question, without asserting that it remains open; no additional characterization problem is invented here. The registry's overall open label and its statement-verification label are supplied metadata, not independent mathematical verification.
\subsection{Short English statement}
The first Erd\H{o}s question asks whether every partial block design whose blocks have size greater than \(c\sqrt n\) admits a set meeting every block in between one and \(C(c)\) points. The second asks for at least \(\exp(c_0\sqrt n\log n)\) distinct realizable nonincreasing block-size sequences for every sufficiently large \(n\). Supplied dated curation reports the counting bound as settled and a structural characterization as unresolved; the intended current open target therefore requires review.
\subsection{Sources}
\begin{itemize}
\item Local UnsolvedMath registry, supplied complete record: permanent ID \(30006389\), source problem code \texttt{catalogue entry}, title ``Two Erd\H{o}s Problems on Partial Block Designs,'' in the Oberwolfach Reports Open Problems collection. The supplied statement attributes the questions to Erd\H{o}s; the proposer field is empty, and the proposed year is recorded as 2025.
\item Supplied report citation: \emph{Combinatorics, Probability and Computing (2026)}, Oberwolfach Reports, DOI 10.4171/owr/2025/42, pp. 2243--2322. \url{https://doi.org/10.4171/owr/2025/42}. The supplied background describes Oberwolfach Reports as documenting mathematical workshops and their open-problem discussions.
\item Supplied literature source: ``Oberwolfach Report 42 2025.'' \url{https://oa.tib.eu/renate/bitstreams/72e281a6-ba33-4ac7-baf7-6fadd1126345/download}.
\item Supplied literature source: ``Alon Blocking Partial Designs and Block-Compatible Sequences.'' \url{https://web.math.princeton.edu/~nalon/PDFS/remark191.pdf}. No further bibliographic details are supplied.
\item Literature assessment date: \texttt{2026-08-22}. This is supplied dated curation derived from the supplied OWR curation, not independently verified. It records status ``open'' and states that the counting lower bound is settled while necessary-and-sufficient conditions for block-compatible sequences remain unresolved. The supplied background expressly calls for independent verification of time-sensitive status and cited claims.
\end{itemize}
\end{document}
